% TOP_PROBLEM: {"id": 1500005, "title": "Algebraic Stories — The $k$-rank of monomials", "category_id": 4, "category": {"id": 4, "name": "algebra", "display_name": "Algebra", "description": "Group theory, ring theory, field theory, and algebraic structures.", "slug": "algebra", "order_index": 4, "created_at": "2026-07-31T15:26:25.671Z"}, "status": "partially_solved"}
% FIRST_POSTED: null
% ENTRY_KIND: statement_only
% Imported from: batch06.tex; UnsolvedMath ID 1500005
% Compile twice: lualatex 1500005.tex
\documentclass[11pt]{article}
\usepackage{fontspec}
\setmainfont{Latin Modern Roman}
\usepackage{amsmath,amssymb,mathtools,unicode-math}
\setmathfont{Latin Modern Math}
\usepackage{xurl,hyperref}
\hypersetup{colorlinks=true,urlcolor=blue,pdftitle={UnsolvedMath Problem 1500005}}
\newcommand{\sref}[2]{\hyperlink{source:#1:#2}{[S#2]}}
\newcommand{\cataloguescope}[1]{\paragraph{Mathematical question.}#1}
\begin{document}
% UnsolvedMath ID 1500005; source catalogue code AMR-014-0005
\setcounter{section}{1500004}
\section{Algebraic Stories — The $k$-rank of monomials}\label{1500005}
{\small\sffamily UnsolvedMath ID 1500005 · Algebra}
\subsection{Definitions and mathematical statement}

\paragraph{Shared notation.}$\mathbb N=\{1,2,\ldots\}$, and $\mathbb Z,\mathbb Q,\mathbb R,\mathbb C$ denote the integers, rationals, reals and complex numbers. Any other conventions are specified locally.


Work over $\mathbb C$. For $n\geq1$, let $S=\mathbb C[x_1,\ldots,x_n]$, with $S_e$ its vector space of homogeneous polynomials of total degree $e$. For $k,d\geq1$ and $f\in S_{kd}$, define
\[\operatorname{rk}_k(f)=\min\left\{s\geq0:f=\sum_{j=1}^{s}g_j^k,\quad g_j\in S_d\right\}.
\]
A general form means one whose coefficient vector lies in a suitable nonempty Zariski open subset of $S_{kd}$. Write $\operatorname{rk}_k^\circ(kd,n)$ for the value on general forms. Let $m=x_1^{a_1}\cdots x_n^{a_n}$ with $a_i\geq1$ and $\sum_i a_i=kd$; variables with exponent zero may be omitted.
\paragraph{Statement recorded in the supplied source.}
Determine $\operatorname{rk}_k(m)$ for every such monomial and every $k\geq3$. \sref{1500005}{2}
\subsection{Short English statement}
Determine the exact number of powers needed to express an arbitrary complex monomial.
\subsection{Sources}
\begin{description}
\item[\hypertarget{source:1500005:1}{S1}] ULAM.AI and the UnsolvedMath Contributors, UnsolvedMath: A Curated Collection of Open Mathematics Problems, record 1500005 (AMR-014-0005). CC BY 4.0. \url{https://huggingface.co/datasets/ulamai/UnsolvedMath}.
\item[\hypertarget{source:1500005:2}{S2}] R. Fröberg, S. Lundqvist, A. Oneto and B. Shapiro, \emph{Algebraic Stories from One and from the Other Pockets}, arXiv:1801.01692 (2018), §1.3, Problem C. \url{https://arxiv.org/abs/1801.01692}.
\end{description}
\label{end:1500005}
\end{document}
